% TOP_PROBLEM: {"problemId": "problem.top500-omission-w045049-potential-density-on-k3-surfaces", "canonicalTitle": "Potential Density of Rational Points on K3 Surfaces", "releaseRank": 483, "primaryDomain": "number_theory_arithmetic_geometry", "primaryDomainLabel": "Number theory & arithmetic geometry", "displayStatus": "open_with_solved_subcases", "releaseStatus": "open_with_solved_subcases"}
% !TeX program = lualatex
% ATTEMPT_STATUS: unresolved
% Model/process: GPT_6_Astra_Ultra; detailed mathematical attempt,
% primary-source check, exact arithmetic diagnostics, adversarial review.
% Prepared 27 September 2026; no general resolution or novelty claimed.
\documentclass[11pt]{article}
\usepackage[margin=25mm]{geometry}
\usepackage{fontspec}
\setmainfont{Latin Modern Roman}
\usepackage{amsmath,amssymb,mathtools}
\usepackage{unicode-math}
\setmathfont{Latin Modern Math}
\usepackage{xurl,hyperref}
\hypersetup{colorlinks=true,urlcolor=blue}
\setcounter{secnumdepth}{0}
\setlength{\parindent}{0pt}
\setlength{\parskip}{5pt}
\setlength{\emergencystretch}{3em}
\begin{document}
\section*{\texorpdfstring{Potential Density of Rational Points on K3 Surfaces}{Potential Density of Rational Points on K3 Surfaces}}\label{483}
\subsection{Definitions and mathematical statement}
Let $K$ be a number field. A K3 surface over $K$ is a smooth projective geometrically integral surface $X$ with trivial canonical line bundle and $H^1(X,\mathcal O_X)=0$. A subset is Zariski dense when it is not contained in any proper Zariski-closed subset. The conjecture is
\[
 \forall X/K\text{ K3}\quad\exists L/K\text{ finite}:\quad
 \overline{X(L)}^{\rm Zar}=X.
\]
The extension $L$ may depend on $X$ and $K$, with no uniform degree bound required. Density over the original field $K$ is not asserted. Infinitely many points on one algebraic curve inside $X$ need not be dense in the surface.
\par\smallskip\textit{Formulation and concept references:} \href{https://arxiv.org/html/2505.13262v2}{[S1]}.

\subsection{Short English statement}
Does every K3 surface over a number field acquire a Zariski-dense set of rational points after some finite field extension?
\subsection{Research attempt}

\paragraph{Scope and source audit (27 September 2026).}
The extension must be chosen once for the surface. Producing each new
point over a new number field, or producing infinitely many points on a
single curve, does not meet the quantifiers. All density assertions below
are geometric: the indicated rational points are dense after passage to
an algebraic closure. This is at least as strong as avoiding proper closed
subsets defined over the chosen number field.

The supplied source [S1], version 2 of 15 October 2025, distinguishes
infinitude from potential density. Its Theorem 1.1 gives an extension of
degree at most 12 with infinitely many points for a degree-two K3 surface.
Its moduli-space density statement concerns a family of surfaces with
infinitely many points; it does not say that the points on every surface
in that family are dense in the surface. Bogomolov--Tschinkel [E1,
Theorem 1.1] prove potential density for elliptic K3 surfaces and those
with an infinite automorphism group. Those are attributed theorems,
not new conclusions of this note.

The route pursued here is propagation within one field. First prove two
elementary density lemmas, then obtain product Kummer examples from
elliptic curves with non-torsion points, and finally examine which geometric
data would let the same mechanism work on an arbitrary K3 surface. An
explicit rational example is worked out far enough to verify non-torsion,
the surface equation, and the need for two independent parameters.
The first obstruction to extending this route is made visible by geometric
Picard number one.

\paragraph{Two elementary density lemmas.}
\textbf{Lemma 483.1 (products of infinite subsets of curves).}
Let $C_1,C_2$ be geometrically integral curves and let
$S_i\subset C_i(L)$ be infinite. Then $S_1\times S_2$ is Zariski dense
in $C_1\times C_2$ over $\overline L$.

\emph{Proof.}
Let $Z$ be a closed subset containing that product. Fix $s_1\in S_1$.
The intersection of $Z$ with $\{s_1\}\times C_2$ contains the infinite
set $\{s_1\}\times S_2$. A proper closed subset of an integral curve
is finite, so the whole fiber is contained in $Z$. Thus for each geometric
point $c_2\in C_2(\overline L)$, the intersection of $Z$ with
$C_1\times\{c_2\}$ contains $S_1\times\{c_2\}$ and is the whole
fiber. Therefore $Z=C_1\times C_2$. This uses closed fibers, not an
identification of the product Zariski topology with a product topology.
\hfill$\square$

\textbf{Lemma 483.2 (infinitely many dense curves in a surface).}
Let $X$ be a geometrically integral surface over $L$. Suppose there are
infinitely many distinct geometrically integral curves $C_i\subset X$,
all defined over $L$, such that $C_i(L)$ is infinite for each $i$.
Then $X(L)$ is Zariski dense.

\emph{Proof.}
The closure of $X(L)$ contains every $C_i$, because its intersection with
$C_i$ contains infinitely many points. A proper closed subset of an
integral surface has only finitely many irreducible curve components,
together with finitely many zero-dimensional components. It cannot
contain infinitely many distinct integral curves.\hfill$\square$

The common field $L$ and distinctness of the curves are essential
hypotheses. Listing infinitely many parametrizations of the same curve
does not satisfy the lemma.

\paragraph{One finite extension creates non-torsion on each elliptic factor.}
Here is a useful consequence of Merel's uniform boundedness theorem [E3].
Only that stated theorem is used; no proof of uniform boundedness is
attempted.

\textbf{Lemma 483.3.}
For any elliptic curve $E/K$ over a number field there is an extension
$L/K$ of degree at most two such that $E(L)$ contains a non-torsion point.

\emph{Proof.}
Choose a short Weierstrass equation $y^2=x^3+Ax+B$ over $K$.
For every integer $t$ outside the finitely many roots of its right side,
choose a square root and obtain
\[
 P_t=(t,\sqrt{t^3+At+B})\in E(L_t),\qquad [L_t:K]\le2.
\]
The points are distinct because their $x$-coordinates are distinct.
Set $D=2[K:\mathbb Q]$. Merel's theorem gives a bound on torsion orders
over number fields of degree at most $D$; if the theorem is stated for
each exact degree, take the maximum of its finitely many bounds for
degrees $1,\ldots,D$. Call the result $B_D$.
If every $P_t$ were torsion, each would lie in
\[
 \bigcup_{1\le m\le B_D}E[m](\overline K).
\]
This union is finite, since $E[m]$ has $m^2$ geometric points in
characteristic zero. That contradicts the infinitely many distinct
$P_t$. Choose one non-torsion $P_t$ and put $L=L_t$.\hfill$\square$

The argument selects a single point and its field; it does not take the
compositum of all quadratic fields. Once selected, every multiple of that
point belongs to the same field. This is exactly the field control needed
for the propagation step.

\paragraph{Product Kummer surfaces: a complete potential-density subcase.}
Let $E_1,E_2$ be elliptic curves defined over $K$, and let
$A=E_1\times E_2$. Write $Y=A/\langle\iota\rangle$, where
$\iota(P,Q)=(-P,-Q)$, and let $X\to Y$ be the minimal resolution.
The Kummer construction says that $X$ is a K3 surface in characteristic
zero [E2, Chapter 1, Example 1.3(iii)]. This is a geometric input, not a
claim that every K3 surface admits this presentation. The involution fixes
the 16 geometric two-torsion points; their quotient singularities are
ordinary double points, and the resolution is an isomorphism away from
their images.

Apply Lemma 483.3 to obtain non-torsion $P_i\in E_i(L_i)$ with
$[L_i:K]\le2$. The compositum $L=L_1L_2$ has degree at most four.
For $S_i=\{nP_i:n\ge1\}$, non-torsion guarantees infinitely many
distinct points, all defined over $L$. Lemma 483.1 proves that
$S_1\times S_2$ is dense in $A$.

The quotient morphism $q:A\to Y$ is finite and dominant. Its image of a
dense subset is dense: if a closed $Z\subset Y$ contained the image,
then $q^{-1}(Z)$ would contain the dense subset and hence all of $A$,
forcing $Z=Y$. None of the positive multiples $nP_i$ is two-torsion, so
the chosen product points avoid the fixed locus. Their images therefore
lie in the open set where $X\to Y$ is an isomorphism, and lift uniquely
to $L$-points of $X$ there. They are dense in that open set and thus in
$X$. We have proved
\begin{equation}\label{ra483:kummer}
 X=\operatorname{Kum}(E_1\times E_2)
 \quad\Longrightarrow\quad
 \exists L/K,\ [L:K]\le4,\quad X(L)\text{ dense}.
\end{equation}
The degree bound belongs only to this product presentation with both
factors already over $K$. It is not asserted for arbitrary K3 surfaces
or arbitrary forms of Kummer surfaces. The method reconstructs a known
special-case mechanism, using the cited Kummer and torsion theorems.

\paragraph{An explicit surface and a non-torsion certificate over $\mathbb Q$.}
Take
\[
 E:y^2=x^3-2,\qquad P=(3,5).
\]
The discriminant of the displayed model is $-1728$, so the reductions at
5 and 7 are good. Direct enumeration gives
\begin{equation}\label{ra483:reductions}
 |E(\mathbb F_5)|=6,\qquad |E(\mathbb F_7)|=7.
\end{equation}
For clarity, at $p=5$ the five right-hand-side values for $x=0,1,2,3,4$
are $3,4,1,0,2$, contributing respectively $0,2,2,1,0$ affine points.
At $p=7$, the only $x$ giving nonzero-square right-hand sides are
$x=3,5,6$, each giving value four and two points. Adding the point at
infinity gives the stated totals.

At a good prime $p$, reduction is injective on torsion of order prime to
$p$. The elementary reason is that the reduction kernel is a formal
group: in its parameter $T\in p\mathbb Z_p$, multiplication by an integer
$m$ prime to $p$ has the form $mT+\sum_{j\ge2}a_jT^j$ with integral
coefficients. If $T\ne0$, the linear term has strictly smaller valuation
than every higher term, so the sum cannot vanish. This proves the needed
prime-to-$p$ injectivity for the rational points used here.
If a rational torsion point had order divisible by a prime $\ell\ne5,7$,
its $\ell$-primary part would inject into both groups in
\eqref{ra483:reductions}, which have coprime orders. If $\ell=5$, use
reduction at 7; if $\ell=7$, use reduction at 5. Every possibility is
excluded. Thus $E(\mathbb Q)$ has trivial torsion. Since $P$ is not the
identity, it is non-torsion.

An affine chart of the quotient of $E\times E$ is
\begin{equation}\label{ra483:surface}
 z^2=(u^3-2)(v^3-2),
 \qquad (u,v,z)=(x(P_1),x(P_2),y(P_1)y(P_2)).
\end{equation}
This equation follows from multiplying the two elliptic equations. More
precisely, the function field invariant under simultaneous sign change
of the two $y$-coordinates is generated by $u,v,z$ with the displayed
relation; the product field is quadratic over it. This is the quotient
chart, not a claim that its naive projective closure is already smooth.
The smooth projective surface under consideration is its Kummer resolution.

For every $m,n\ge1$, the point
\[
 Q_{m,n}=(x(mP),x(nP),y(mP)y(nP))
\]
is rational and lies in this chart. Positive multiples of $P$ are neither
the identity nor two-torsion, so the coordinates are finite and avoid the
singular quotient points. The preceding product argument proves that their
lifts are Zariski dense in the K3 surface. This explicit example already
has density over $\mathbb Q$; no extension is needed.

The first doubling is a useful exact check. Its tangent slope is $27/10$,
so
\[
 2P=\left(\frac{129}{100},-\frac{383}{1000}\right),\qquad
 Q_{1,2}=\left(3,\frac{129}{100},-\frac{383}{200}\right).
\]
Both equations are verified by direct rational substitution. Keeping only
$Q_{n,n}$ would fail to prove density: that sequence satisfies $u=v$ and
$z=u^3-2$, and therefore lies on one curve of the surface. This is an
explicit demonstration of why the two independent parameters in the
successful construction are doing essential work.

\paragraph{An elliptic-fibration propagation lemma.}
Suppose a smooth projective surface $X/L$ has an elliptic fibration
$\pi:X\to\mathbb P^1_L$ with zero section and a section $P$ of infinite
order on the generic elliptic fiber over $L(t)$. This is a precise
additional hypothesis, not a property claimed for every K3 surface.

For each positive integer $n$, the generic point $nP$ gives a rational
section. It extends over the whole base: a rational map from a smooth
complete curve to a proper variety extends at every missing point by the
valuative criterion of properness. Its composite with $\pi$ agrees with
the identity on the generic point and hence everywhere. Thus the section
image $C_n$ is a curve isomorphic to $\mathbb P^1_L$. If $C_n=C_m$, the
two sections agree on the generic fiber, since a section has a unique point
over that generic base point. This gives $(n-m)P=0$, so $n=m$.

All $C_n$ are therefore distinct curves defined over the single field
$L$, and each has infinitely many $L$-points. Lemma 483.2 proves density
of $X(L)$. Notice what was not required: no assertion that almost every
specialization of $P$ is non-torsion, and no uniform specialization theorem.
Density is propagated across distinct section curves instead of proved
separately on every fiber.

If the fibration and non-torsion section are initially given over
$\overline K$, their equations and morphisms involve finitely many
algebraic coefficients. Adjoining these finitely many coefficients gives
one finite extension $L/K$ over which the construction is defined. The
group law then defines all multiples over that same field. This finite
descent of a fixed set of generating data is valid; descending an arbitrary
infinite list of unrelated curves is not the same argument.

The theorem of [E1] covers elliptic K3 surfaces more generally, including
cases in which this particular positive-rank section hypothesis is not
available. The elementary lemma is a convenient sufficient subcase, not
a replacement proof for the entire elliptic theorem.

\paragraph{A second propagation mechanism, with its needed hypothesis.}
Suppose $\sigma:X\to X$ is an automorphism defined over $L$, and
$C\subset X$ is a geometrically integral curve defined over $L$ with
infinitely many $L$-points. If the curves $\sigma^n(C)$ are infinitely
many distinct curves, each has infinitely many $L$-points, and Lemma 483.2
again gives density over this same field.

Infinite order of $\sigma$ alone does not verify the distinctness condition
for a chosen $C$: an infinite-order automorphism can stabilize a curve.
Similarly, infinitely many iterates of one point could all remain on one
invariant curve. The successful input is an infinite orbit of the curve
itself, together with arithmetic density on that curve. An arbitrary K3
surface is not supplied here with such a curve-automorphism pair.

\paragraph{Why extending the field does not manufacture every needed fibration.}
There is a simple geometric obstruction to extending the preceding
fibration proof to all K3 surfaces. Suppose the geometric Neron--Severi
space has rank one, generated over $\mathbb Q$ by an ample class $H$.
If a fibration to a curve existed over $\overline K$, the class $F$ of a
fiber would satisfy $F^2=0$: distinct fibers are disjoint and numerically
equivalent. Also $H\cdot F>0$, so $F\ne0$. But rank one forces
$F=aH$ for a nonzero rational $a$, giving
\[
 F^2=a^2H^2>0,
\]
a contradiction. Thus a geometrically rank-one K3 surface does not acquire
an elliptic fibration by passing to a finite extension. Such an extension
can make existing geometric divisor classes rational; it cannot change
the geometric Neron--Severi space itself.

The original source's rank-one examples with infinitely many points are
therefore relevant stress tests. Those infinitely many points may lie on
the same curve constructed in its argument. Neither their infinitude nor
the density of the corresponding surface family in a moduli space supplies
surface-wise potential density. This distinction is exactly what the
catalogue formulation retains.

\paragraph{The finite-extension quantifier cannot be exchanged.}
For any variety of finite type over a number field, every geometric point
has coordinates in some finite extension. Therefore
\[
 X(\overline K)=\bigcup_{L/K\ \mathrm{finite}}X(L).
\]
The left side is dense for elementary geometric reasons. This observation
does not produce a single $L$ for which $X(L)$ is dense. The elementary
field example $\mathbb Q(\sqrt p)$ for distinct primes $p$ already shows
that infinitely many data sets of bounded degree need not fit inside one
finite extension: a fixed number field has only finitely many intermediate
quadratic fields, as follows by passing to its finite Galois closure.

In the Kummer proof, the infinite set came from two points after one
finite compositum. In the section proof, it came from one section using
group operations over one field. These are actual mechanisms preserving
the field. No corresponding finite generating data for all the necessary
curves or points has been produced for an arbitrary K3 surface.

\paragraph{Verification and outcome.}
The companion script enumerates both finite reduction groups in
\eqref{ra483:reductions}, calculates the first twelve positive multiples
of $P$ with exact rational group-law arithmetic, and checks that they are
distinct and satisfy $y^2=x^3-2$. It verifies all 144 points $Q_{m,n}$
obtained from those multiples in \eqref{ra483:surface}, together with the
explicit $Q_{1,2}$ above. This is a finite check of the formulas; infinitude
and density come from the non-torsion and product proofs, not the table.

The singular affine quotient is always distinguished from its smooth
projective resolution. Only points in the isomorphism locus are lifted,
so no unwarranted rational point on an exceptional conic is assumed.
The sections are proved distinct on the generic fiber, and the common
field is fixed before their infinitely many multiples are formed.

\textbf{Outcome: UNRESOLVED, with complete propagation subcases.}
The attempt proves potential density for products of elliptic curves
passed to their Kummer surfaces, with an explicit rational dense example,
and proves two conditional surface propagation lemmas. It does not cover
an arbitrary K3 surface, especially geometric Picard number one. The first
remaining gap is to construct enough arithmetically dense curves or other
point-producing data over one finite extension when these group and
fibration structures are absent. Neither a general construction nor a
K3 counterexample to potential density is obtained.

\subsection{Sources}
\begin{itemize}
\item[S1]J. Martínez-Marín, Rational points on K3 surfaces of degree 2, version 2 (2025), Introduction.
\url{https://arxiv.org/html/2505.13262v2}.
\textit{Formulation source.}
\item[E1]Fedor A. Bogomolov and Yuri Tschinkel,
\emph{Density of rational points on elliptic K3 surfaces},
Asian Journal of Mathematics 4 (2000), no.~2, 351--368;
arXiv:math/9902092, Theorem 1.1.
\url{https://arxiv.org/pdf/math/9902092}.
\textit{Primary theorem inspected 27 September 2026; the general
elliptic/infinite-automorphism theorem is attributed, not reproved.}
\item[E2]Daniel Huybrechts, \emph{Lectures on K3 Surfaces},
author's draft, Chapter 1, Example 1.3(iii), pp.~8--9
(Kummer construction), and Section 2 (intersection pairing).
\url{https://www.math.uni-bonn.de/people/huybrech/K3Global.pdf}.
\textit{Relevant pages directly inspected 27 September 2026;
draft pagination and numbering are those of this linked file.}
\item[E3]Loic Merel, \emph{Bornes pour la torsion des courbes elliptiques
sur les corps de nombres}, Inventiones Mathematicae 124 (1996), 437--449,
opening uniform-boundedness corollary.
\url{https://perso.imj-prg.fr/wp-content/uploads/merel-pub/torsion.pdf}.
\textit{Primary author PDF inspected 27 September 2026; used only for
bounded-degree torsion, with the quadratic-extension consequence proved here.}
\end{itemize}
\end{document}
